\chapter[1981 --- Sperry et al.]{1981: Roger W. Sperry, David H. Hubel, Torsten N. Wiesel}
\label{ch:1981_Sperry}

\section*{Main Contribution}

\textit{Discovered the functional specialization of the cerebral hemispheres, explaining how the left and right brain process different types of information.}

\section{Official Citation}

for his discoveries concerning the functional specialization of the cerebral hemispheres

\section[Roger W. Sperry: Discoveries]{Roger W. Sperry: Discoveries Concerning the Functional Specialization of the Two Hemispheres of the Brain}

\subsection{Background and Historical Context}

The 1981 Nobel Prize in Physiology or Medicine was divided: one half was awarded to Roger W. Sperry ``for his discoveries concerning the functional specialization of the two hemispheres of the brain,'' and the other half was awarded jointly to David H. Hubel and Torsten N. Wiesel ``for their discoveries concerning information processing in the visual system.'' Together, their work revealed fundamental principles of brain organization---how the brain processes visual information and how the two cerebral hemispheres specialize for different cognitive functions.

For centuries, the human brain has been an object of fascination and mystery. The two cerebral hemispheres, connected by the corpus callosum, were known to be roughly symmetrical in structure, and it was generally assumed that they performed similar or identical functions. This assumption was reinforced by the prevailing view, articulated by the Nobel laureate John Eccles and others, that the left hemisphere was dominant for language and intellectual function, while the right hemisphere was essentially a ``minor'' or ``ancient'' hemisphere that played little or no role in higher cognitive function.

Roger Wolcott Sperry was born in Hartford, Connecticut, in 1913, and trained in psychology at Oberlin College and at the University of Chicago, where he earned his Ph.D. in 1941. After a postdoctoral fellowship with Karl Lashley at Harvard, Sperry joined the faculty of the University of Chicago and later moved to the California Institute of Technology (Caltech), where he conducted his Nobel Prize-winning work on the split brain.

\subsection{The Discovery/Contribution}

Sperry's investigations of the functional specialization of the cerebral hemispheres began with a series of experiments on the corpus callosum---the thick bundle of nerve fibers that connects the two hemispheres of the brain. In the 1950s and 1960s, Sperry and his colleagues at Caltech conducted experiments on cats and monkeys in which the corpus callosum had been surgically severed (a procedure known as commissurotomy) to prevent the spread of epileptic seizures between hemispheres. They showed that after commissurotomy, the two hemispheres became functionally independent---each hemisphere could learn and remember things that the other hemisphere knew nothing about.

Sperry's primary insights, however, came from his studies of human patients who had undergone commissurotomy for the treatment of severe epilepsy. Working with his colleague Michael Gazzaniga in the late 1960s, Sperry developed a series of ingenious behavioral tests to assess the cognitive abilities of each hemisphere independently. In these tests, visual stimuli were presented to one hemisphere or the other (by briefly flashing images to the left or right visual field), and the patient's responses were recorded.

The results were notable. Sperry and Gazzaniga showed that the two hemispheres were not equivalent in their cognitive abilities. The left hemisphere was dominant for language, speech, and analytical reasoning---it could name objects, describe relationships, and perform mathematical calculations. The right hemisphere, by contrast, was superior in spatial perception, face recognition, musical ability, and the processing of emotional and social information---but it was largely unable to produce speech or to understand complex linguistic structures.

These findings overturned the long-held assumption that the right hemisphere was cognitively inferior to the left. Sperry showed that the right hemisphere was, in many respects, more capable than the left---for example, it was better at recognizing faces, at performing spatial tasks, and at understanding the emotional content of speech. However, the two hemispheres were largely unable to communicate with each other after the corpus callosum had been severed, and each hemisphere was essentially ``unaware'' of the other hemisphere's experiences and knowledge.

Sperry's work had significant implications for our understanding of consciousness, lateralization, and the organization of the brain. It demonstrated that the brain is not a single, unified organ, but rather a collection of specialized modules that perform different functions. It also showed that consciousness itself might be divided---that a single individual could harbor two separate streams of awareness, each with its own perceptions, memories, and intentions.

\section[David H. Hubel and Torsten N.]{David H. Hubel and Torsten N. Wiesel: Discoveries Concerning Information Processing in the Visual System}

\subsection{Background and Historical Context}

The visual system has long been a favored model for studying how the brain processes sensory information. The retina, the optic nerve, and the visual cortex---the area of the brain that receives and processes visual information---are well-defined and accessible structures, and the visual cortex is organized in a systematic, topographic manner that lends itself to experimental analysis.

David Hunter Hubel was born in Windsor, Ontario, Canada, in 1926, and trained in medicine at McGill University. After completing his medical training, Hubel joined the neurophysiology laboratory of Stephen Kuffler at Johns Hopkins University, where he began recording from individual neurons in the visual cortex of cats. Torsten Nils Wiesel was born in Uppsala, Sweden, in 1924, and also trained in medicine at the Karolinska Institute. He joined Kuffler's laboratory at Johns Hopkins in 1955, and he and Hubel began a collaboration that would last for more than two decades and produce some of a major discoveries in the history of neuroscience.

\subsection{The Discovery/Contribution}

Hubel and Wiesel's work on the visual system began in the late 1950s, when they developed techniques for recording the electrical activity of individual neurons in the visual cortex of anesthetized cats. Their approach was to project patterns of light and shadow onto a screen in front of the cat's eyes and to record the responses of individual cortical neurons to these stimuli.

The first major discovery came in 1959, when Hubel and Wiesel found that neurons in the primary visual cortex (area 17, or V1) responded not to spots of light (as retinal ganglion cells do) but to elongated edges and bars of light at specific orientations. Some neurons responded strongly to a vertical bar of light, others to a horizontal bar, and still others to bars at various oblique angles. Hubel and Wiesel called these neurons ``simple cells'' because their responses could be predicted by a simple mathematical model based on the arrangement of excitatory and inhibitory regions in the receptive field.

Further investigation revealed a hierarchy of increasingly complex feature-detecting neurons in the visual cortex. ``Complex cells'' responded to bars of light at specific orientations regardless of their exact position within the receptive field, and ``hypercomplex cells'' (later called ``end-stopped cells'') responded to bars of light of specific lengths. Hubel and Wiesel proposed that the visual cortex is organized as a hierarchical feature-detection system, in which simple features detected by neurons at one level are combined to form more complex representations at higher levels.

Hubel and Wiesel also discovered that the visual cortex is organized into columns---vertical arrays of neurons that share similar orientation preferences. Neurons in the same column respond to bars of light at the same orientation, while neurons in adjacent columns respond to slightly different orientations. This columnar organization suggested that the visual cortex is organized as a modular system in which all possible orientations are represented in a systematic, orderly fashion.

In addition to the orientation columns, Hubel and Wiesel discovered ocular dominance columns---alternating stripes of cortex that receive input preferentially from one eye or the other. In the primary visual cortex, neurons in ocular dominance columns respond preferentially to input from the left eye or the right eye, and the inputs from the two eyes are combined at higher levels of the visual hierarchy to produce binocular vision and depth perception.

Hubel and Wiesel's primary contribution to understanding visual processing, however, was their demonstration that the visual cortex is not a static, hard-wired structure but rather a dynamic system that is shaped by experience during a critical period of early development. In a series of experiments conducted in the 1960s and 1970s, Hubel and Wiesel showed that if one eye was sutured closed during the critical period of visual development (the first few months of life in kittens), the neurons in the visual cortex that would normally have received input from the deprived eye became responsive to input from the open eye instead. This ``ocular dominance plasticity'' demonstrated that the connectivity of the visual cortex is not genetically predetermined but is shaped by visual experience during a critical window of development.

The critical period concept---the idea that there is a limited time window during which the brain is particularly sensitive to environmental influences---has proved to be of significant importance for understanding brain development, learning, and recovery from brain injury. It has implications for the treatment of amblyopia (lazy eye), for the rehabilitation of stroke patients, and for our understanding of how early childhood experiences shape brain function.

\section{Impact on Modern Medicine}

The discoveries of Sperry, Hubel, and Wiesel have had a significant impact on modern medicine and neuroscience. Sperry's work on the lateralization of brain function has influenced the diagnosis and treatment of neurological disorders, including epilepsy, stroke, and traumatic brain injury. Understanding the specialized functions of the two hemispheres has enabled clinicians to predict the functional consequences of lesions in different brain regions and to develop targeted rehabilitation strategies. Sperry's work has also informed the development of neuroimaging techniques, such as functional MRI (fMRI), which can visualize the activation of different brain regions during cognitive tasks.

Hubel and Wiesel's work on the visual system has had a particularly broad impact. Their discoveries about the organization of the visual cortex have provided the foundation for our understanding of visual processing, visual perception, and visual disorders. The concept of critical periods has been applied to the treatment of amblyopia, in which early intervention (patching the stronger eye to force the weaker eye to develop) can restore normal vision if treatment is initiated during the critical period. The critical period concept has also influenced our understanding of language development, in which early exposure to language is essential for normal language acquisition.

Hubel and Wiesel's work has also influenced the development of artificial intelligence and computer vision. The hierarchical feature-detection system that they described in the visual cortex has inspired the design of convolutional neural networks (CNNs), which are now the most powerful algorithms for image recognition and computer vision. CNNs, which are used in applications ranging from medical image analysis to autonomous driving, are directly inspired by the hierarchical organization of the visual cortex that Hubel and Wiesel elucidated.

In the field of ophthalmology, Hubel and Wiesel's work has provided insights into the mechanisms of visual deprivation and the importance of early visual experience for normal visual development. Their findings have informed the design of screening programs for amblyopia and other visual disorders in children, and they have contributed to the development of new treatments for visual impairment.

\section{Legacy}

Roger Sperry continued his work at Caltech until his retirement, exploring the relationship between brain lateralization and consciousness, creativity, and human values. He became an influential advocate for the scientific study of consciousness and for interdisciplinary approaches to understanding the mind. Sperry died in 1994, but his insights into the lateralization of brain function continue to influence neuroscience and psychology.

David Hubel and Torsten Wiesel continued their collaboration at Harvard Medical School, where they made further contributions to our understanding of the visual cortex and the mechanisms of visual development. Hubel died in 2013, and Wiesel continued to be active in neuroscience research and advocacy until his death in 2024. Wiesel was also a prominent advocate for science education and for the protection of children in conflict zones.

The work of the three laureates of 1981 revealed fundamental principles of brain organization that continue to shape our understanding of how the brain processes information, how it is shaped by experience, and how it gives rise to the richness of human perception and cognition. Their discoveries have saved and improved countless lives through better diagnosis and treatment of neurological and visual disorders, and they continue to inspire new generations of scientists and clinicians.


\section{Modern Developments}

Sperry, Hubel, and Wiesel's work on brain lateralization and visual processing has led to modern neuroscience. Understanding of visual cortex organization has informed development of visual prostheses. Understanding of critical periods in development has led to treatments for amblyopia (lazy eye). Understanding of cortical plasticity has informed rehabilitation after stroke and brain injury. Brain-computer interfaces decode neural signals from visual cortex to restore vision.
